\chapter*{读者指南与记号约定}
\phantomsection
\addcontentsline{toc}{chapter}{读者指南与记号约定}

初读时，每节先弄清正在研究的物体、观察方式和要算的量。遇到新符号，回到它第一次出现的等式，找出它由哪些已知量定义。遇到较长的推导，可以先用一句话说出每一步在做什么，例如投影、积分、用初值定常数，再自己写出中间的等式。

例题可以先盖住解答，试着画图并写第一条方程。章节末尾先做基础题，再做综合题。拓展和计算部分可以在掌握核心模型之后回读。下表可用来查找跨章重复使用的符号。

\medskip
\begin{center}
\begin{tabular}{@{}p{3.3cm}p{10.7cm}@{}}
\toprule
记号 & 含义与使用条件\\
\midrule
$\br,\bv,\ba$ & 位置、速度和加速度；比较分量前先注明参考系和基矢。\\
$T,V,E_{\rm mech}$ & 经典动能、位置势能、机械能；在相应模型中 $E_{\rm mech}=T+V$。\\
$L,S,E_L,H$ & 拉格朗日量、作用量、能量函数、哈密顿量。能否将 $E_L$ 或 $H$ 认作机械能，须核对模型与变量。\\
$\bp,\boldsymbol{\pi}$ & 电磁耦合中的广义动量（第6章称正则动量）和机械动量；自由质点两者相同。\\
$e,e_r$ & 电磁模型的带符号电荷、碰撞模型的恢复系数；轨道例中的 $e$ 也可局部表示偏心率。\\
$\theta,\varphi$ & 球坐标中为极角、方位角；第1章的平面极角 $\theta$ 在第3、6章的平面中心力模型中也记作 $\varphi$。\\
$m,M,E,K$ & 第7章分别为粒子不变质量、系统不变质量、粒子总能量与动能；$K=E-mc^2$。本书取时空度规号差 $(+,-,-,-)$。\\
\bottomrule
\end{tabular}
\end{center}

\textbf{怎样核对结果。}
先读题设与假设，再看推导和量纲；若使用守恒量，先说明系统边界。核心结论框归纳公式及条件，物理图像框解释机制，数学注记说明判据，进一步探索框给出延伸主题。各章习题按“章号.题号”编号，标出基础、综合或拓展；解答使用同一题号。计算题可用随书脚本检查残差、守恒或步长误差。
